The three tests rest on four bodies of work: conformal prediction and the law of its coverage, tests of a proportion and of interval forecasts, the detection and the evaluation of dataset shift, and situation testing in fairness. A fifth body, the methodology of simulation studies, supplies the standard under which the error rate of a rule is measured. This section characterises each, states what it assumes and what it gives, and says where the corrections stand.

\subsection{Conformal prediction and the law of its coverage}

Conformal prediction builds, from a score of non-conformity and a calibration sample, a prediction set whose coverage is guaranteed under exchangeability \citep{vovk2022}. The inductive variant splits the data into a proper training set and a calibration set, so that the score quantile is computed once on the calibration rows \citep{papadopoulos2002}. For regression this split-conformal construction has finite-sample marginal coverage of at least $1-\alpha$, and an upper bound that exceeds the target by a term of the order of the reciprocal of the calibration size \citep{lei2018}. Conformalised quantile regression adapts the width of the set to the covariates while keeping the same guarantee \citep{romano2019}. The guarantee is marginal over the calibration sample and the test point. Conditional on one calibration sample, the coverage of the set is a random variable: Vovk showed that, under exchangeability with continuous scores, it has a Beta law whose parameters are the rank of the quantile and the calibration size \citep{vovk2012}. Angelopoulos and Bates restate that law and, in an appendix, give the Beta-Binomial law of the empirical coverage on a validation window of $m$ points, its mean and its variance, as a way to debug an implementation: if the observed coverage is far from what the law allows, the implementation is suspect \citep{angelopoulos2021}.

Two lines of work relax exchangeability. Weighted conformal prediction restores the guarantee under a known covariate shift by weighting the calibration scores with the likelihood ratio of the two covariate laws \citep{tibshirani2019}, and a bound on the coverage gap under a general violation of exchangeability depends on how far the data are from exchangeable and on the weights chosen \citep{barber2023}. Adaptive conformal inference updates the target level online, so that the long-run average coverage holds under an arbitrary shift \citep{gibbs2021}. These works make the set valid under a shift. A validation test asks the opposite question: whether the shift has broken the set, on a window of rows with labels, and how often that question is answered wrongly when nothing has changed. The law of the covered count given the calibration size is published; a test built on it, its exact size, and the measurement of the error of the binomial test that the validation practice names are not stated in the works above, and they are what the first correction of this manuscript supplies.

\subsection{Tests of a proportion and of interval forecasts}

The oldest coverage tests are tests of a proportion. The exact binomial interval inverts the binomial distribution and is conservative by construction \citep{clopper1934}; the Wald interval is simple and misbehaves in small samples, for which the score interval of Wilson and its relatives are preferred \citep{wilson1927,brown2001}. In risk management the same logic gives the proportion-of-failures test of value-at-risk exceedances \citep{kupiec1995} and the tests of unconditional coverage, independence and conditional coverage of interval forecasts \citep{christoffersen1998}. All of them treat the nominal rate as known. A conformal set differs in one respect: its realised coverage is itself estimated from a finite calibration sample, so the nominal rate is known only up to a Beta fluctuation. The distribution that results from mixing a binomial over a Beta rate is the Beta-Binomial distribution, whose probability function and moments are standard \citep{johnson2005}. In the validation of probability-of-default estimates the same family of error has been described: a binomial test of a calibration level assumes that the defaults are independent, and when the defaults are correlated, the test rejects a correct model too often \citep{tasche2008}. The studies of the Basel Committee on the validation of internal rating systems review these tests and the quantitative tools around them \citep{bcbswp14}. The corrected coverage test belongs to this family; it differs in the source of the extra variance, which is the calibration quantile, and in its use of the exact law instead of a normal approximation.

\subsection{Dataset shift and the evaluation of a model under shift}

Dataset shift has a taxonomy, in which covariate shift, prior-probability shift and concept shift are the main forms \citep{morenotorres2012}. Detecting it without labels combines a dimension reduction with a two-sample test; an empirical comparison of such detectors runs combinations of a dimension reduction and a two-sample test on shifts of several kinds \citep{rabanser2019}. The two-sample test based on a classifier, whose accuracy on held-out data of two samples is compared with chance, is the statistical core of the resemblance score used here to rank the evaluation rows by how little they resemble the development rows \citep{lopezpaz2017}. Label shift has its own estimator, which uses a black-box predictor and its confusion matrix to estimate the change in the label proportions \citep{lipton2018}. These methods answer whether a shift exists. The resilience test of the PiML toolbox asks a different question: how much performance the model loses on the subsample of the evaluation rows that is the worst by some criterion, and it compares the loss with a threshold \citep{sudjianto2023}. The worst subsample is a region of the feature space, and the area under the ROC curve on a restricted range differs from the area on the full range even for a model that has not changed; the interpretation of the area as a probability of a correct ranking, and its sampling variance, are standard results \citep{hanley1982,delong1988}. A threshold on the raw drop therefore has a size that no work on the test states.

Permutation and Monte Carlo tests give an exact size when the observed statistic and the simulated ones are exchangeable \citep{davison1997}. When the permutations are drawn at random, the p-value $(1+b)/(1+R)$, with $b$ the number of draws at least as extreme as the observation and $R$ the number of draws, is a valid p-value that is never zero \citep{phipson2010}. The second correction of this manuscript is a Monte Carlo test of this kind, whose null draws are pseudo-evaluation subsets of the development rows.

\subsection{Situation testing}

Situation testing compares the treatment of an individual of a protected group with that of the individuals most similar to her in the control and in the test group. The nearest-neighbour implementation selects, for each complainant, the $k$ most similar members of each group and compares their rates of negative decisions \citep{luong2011}. The counterfactual extension builds the comparison group around a counterfactual version of the complainant and gives a one-sided interval on the difference in rates, with a per-complainant significance \citep{alvarez2025}. Neither work addresses what share of significant complainants should trigger a flag at the level of a portfolio, because both treat the individual case. The light of the engine turns the per-complainant results into a portfolio-level indicator, the share of significant complainants, and a cut point on that share is a convention. The measurement of the size of that convention is the third correction of this manuscript. The data on which the committed fairness case runs, the public mortgage-decision data of the United States, are the ones the cut meets in practice \citep{hmda2018}.

\subsection{The methodology of a simulation study}

The error rate of a rule is the result of a simulation study, and the quality of such a study depends on its design: the aims, the data-generating mechanisms, the estimands, the methods and the performance measures, with the Monte Carlo error of each measure reported \citep{morris2019}. The statistical literature on the design of simulation studies adds the guidance on the number of replications and on the reporting of the uncertainty of a rate \citep{burton2006}. Both works state that a simulated rate without its standard error cannot be compared with another. The standard error shrinks only with the square root of the number of replications, so where the null law is known in closed form the exact size, which has no sampling error at all, is the better reference, and the simulation then serves as a check of the code. The protocol of this manuscript follows them: every rate carries the standard error $\sqrt{\hat r(1-\hat r)/R}$ of its $R$ replications, the exact size is computed by summation where the null law is known, and the plants that generate the data have a known truth. The Polish company bankruptcy data, on which the committed coverage evidence rests, are described in the work that introduced them \citep{zieba2016}.

\subsection{Where the corrections stand}

The published methods supply the guarantee (conformal prediction), the law of the coverage, the form of a test of a proportion, a taxonomy and detectors of shift, a worst-subsample evaluation, a Monte Carlo test, a nearest-neighbour test of discrimination, and the methodology of a simulation study. They do not supply, for the three tests of this manuscript, the error rate of the rule a validation team would write from them. The contribution of the manuscript is the measurement of that error rate for the rule and for the correction, and the three corrections themselves, each of which composes known parts: the Beta-Binomial law of the covered count, a permutation reference built from development rows, and a cut point set from a measured size and power. The comparison with the published methods is made by running them: the binomial rule, the fixed cut on the drop and the rule that treats the complainants as independent are the rules that a validation team would write from the literature, and the product runs each of them on the same replications as its correction.
